\section{Задание}
\subsection{Цель работы}
Целью лабораторной работы является построение доверительных интервалов для 
математического ожидания и дисперсии нормальной случайной величины.

\subsection{Содержание работы}
\begin{itemize}
    \item Для выборки объема $n$ из генеральной совокупности $X$ реализовать в
        виде программы на ЭВМ:
        \begin{itemize}
            \item вычисление точечных оценок $\hat \mu(\vec x_n)$ и $S^2(\vec
                x_n)$ математического ожидания $MX$ и дисперсии
                $DX$ соответственно;
            \item вычисление нижней и верхней границ $\underline \mu (\vec
                x_n)$, $\overline \mu (\vec x_n)$ для
                $\gamma$-доверительного интервала для математического
                ожидания $MX$;
            \item вычисление нижней и верхней границ $\underline \sigma^2 (\vec
                x_n)$, $\overline \sigma^2 (\vec x_n)$ для
                $\gamma$-доверительного интервала для дисперсии
                $DX$;
        \end{itemize}
    \item вычислить $\hat \mu(\vec x_n)$ и $S^2(\vec x_n)$ для выборки из
        индивидуального варианта;
    \item для заданного пользователем уровня доверия $\gamma$ и $N$ --- объема
        выборки индивидуального варианта:
        \begin{itemize}
            \item на координатной плоскости $Oyn$ построить прямую \mbox{$y=\hat \mu
                (\vec x_N)$}, также графики функций $y=\hat \mu (\vec x_n)$,
               $y=\underline \mu (\vec x_n)$ и $y=\overline \mu (\vec x_n)$
               как функций объема $n$ выборки, где $n$ изменяется от $1$
               до $N$;
            \item на координатной плоскости $Ozn$ построить прямую
                \mbox{$z=S^2 (\vec x_N)$}, также графики функций $z=S^2 (\vec
                x_n)$, $z=\underline \sigma^2 (\vec x_n)$ и \mbox{$z=\overline
                \sigma^2 (\vec x_n)$} как функций объема $n$ выборки, где $n$
                изменяется от $1$ до $N$.
        \end{itemize}

\end{itemize}